% !TEX root = ../../main.tex
% Part II (Nguyen Huu Cuong) - merged from the research report (report_latex_v3), see MERGE_NOTES.md.
% Style pass 2026-09-30: em dashes and rhetorical sentences removed; technical content, numbers and references unchanged.
\chapter{Attribute Index and Multi-Agent Reasoning}
\label{chapter:r-reasoning}

\textit{This chapter describes the attribute index built offline with a vision-language model, and the multi-agent funnel that re-ranks the candidates of CLIP retrieval from that evidence.}

\section{The Attribute Index}

\subsection{Schema}
Each described crop carries the 18 attributes of Table~\ref{tab:r-attrs}. The
schema is deliberately centred on garments, since these are what a witness
usually describes. \at{carrying\_canon} is list-valued (a person can carry a
backpack \emph{and} a handbag); every other attribute is single-valued per crop,
although the same attribute read across a track often produces several values,
which is the signal the reasoning stages use.

\begin{table}[H]
\centering
\small
\begin{tabular}{@{}lll@{}}
\toprule
\textbf{Attribute} & \textbf{Region} & \textbf{Meaning}\\
\midrule
\at{gender}             & whole & apparent gender\\
\at{age\_group}         & head  & apparent age band\\
\at{build}              & whole & apparent build\\
\at{height\_class}      & whole & apparent height, words only\\
\at{hair\_color}        & head  & dominant hair colour\\
\at{hair\_style}        & head  & visible hairstyle\\
\at{upper\_inner\_type} & upper & the top worn under any outer layer\\
\at{upper\_inner\_color}& upper & the colour of that top\\
\at{upper\_outer\_type} & upper & the outermost upper garment (jacket, coat)\\
\at{upper\_outer\_color}& upper & the colour of that outer garment\\
\at{sleeves}            & upper & sleeve length\\
\at{pattern}            & upper & visible pattern on the upper body\\
\at{lower\_type}        & lower & lower-body garment\\
\at{lower\_color}       & lower & its colour\\
\at{lower\_length}      & lower & how far it reaches\\
\at{shoes\_type}        & feet  & footwear type\\
\at{shoes\_color}       & feet  & footwear colour\\
\at{carrying\_canon}    & whole & non-clothing objects carried (LIST-valued)\\
\bottomrule
\end{tabular}
\caption{The attribute schema, and the visual band each attribute is judged
from. The region column is a manual routing table: a wrong entry is invisible at
run time, because the evidence line still reads correctly but quotes the
sharpness of the wrong part of the body.}
\label{tab:r-attrs}
\end{table}

\subsection{Visual context}
Attributes alone cannot explain a disagreement. If one frame reads the top as red
and another as black, the useful question is whether the two frames were taken
under conditions that make this confusion likely. Every crop therefore also
carries numeric quality statistics, computed per band:

\begin{center}
\begin{tabular}{@{}ll@{}}
\toprule
\textbf{Band} & \textbf{Extent of the crop}\\
\midrule
\at{whole} & the entire crop\\
\at{head}  & the first 15\%\\
\at{upper} & the next 30\%\\
\at{lower} & the next 40\%\\
\at{feet}  & the last 15\%\\
\bottomrule
\end{tabular}
\end{center}

\noindent
Each band carries four numbers: sharpness, contrast, brightness and saturation.
Because \at{upper\_inner\_color} is routed to the \at{upper} band, a note about a
colour disagreement can cite the saturation of the upper body in the frames that
read red and in those that read black, giving a mechanism with its figure rather
than a guess.

\subsection{Three different counts}
The same attribute value can be counted in three ways. Confusing them is the
easiest mistake in this design, so the implementation keeps them as three
separate functions:

\begin{itemize}
\item The \emph{global prior} is the value's share over the \emph{whole index},
      with lists \emph{split}: a reading of \at{[red, white]} plus a reading of
      \at{red} gives red 66\% and white 33\% over three observations. This is
      the base rate against which rarity is judged.
\item The \emph{track distribution} is the set of readings of one track, kept
      \emph{whole}: the same two readings give ``red+white 50\%, red 50\%''.
      This is what the evidence reader quotes.
\item \emph{Membership} tells whether a query value \at{red} matches a frame that
      read \at{[red, white]}, which it does.
\end{itemize}

The local (pool) prior adds one more rule: it is computed over
\emph{deduplicated tracks}, not over pool crops. Two pool crops from the same
five-frame track give a five-frame statistic, not a ten-frame one. Without this
rule, a person retrieved many times would dominate the pool statistic simply
because the retriever returned that person several times.

\subsection{The image-quality problem, and why voting does not work}
\label{sec:r-noise}
Table~\ref{tab:r-noise} shows one real track, described over ten frames. It is
the motivating example for Section~\ref{sec:r-reasoning}.

\begin{table}[H]
\centering
\small
\begin{tabular}{@{}ll@{}}
\toprule
\textbf{Attribute} & \textbf{Across the 10 described frames}\\
\midrule
\at{gender}             & female 80\%, male 20\%\\
\at{age\_group}         & young adult 100\%\\
\at{build}              & slim 100\%\\
\at{height\_class}      & average height 80\%, average 20\%\\
\at{hair\_color}        & black 40\%, brown 30\%, blonde 30\% (split)\\
\at{hair\_style}        & short 50\%, straight 30\%, long+straight 20\% (split)\\
\at{upper\_outer\_type} & hoodie 20\%, not visible 80\% (split)\\
\at{upper\_outer\_color}& black 20\%, not visible 80\% (split)\\
\at{upper\_inner\_type} & sweater 70\%, tank top 10\%, not visible 20\%\\
\at{upper\_inner\_color}& gray 40\%, black 40\%, not visible 20\%
                          (tied, the aggregate is arbitrary)\\
\bottomrule
\end{tabular}
\caption{VLM noise on one track. A majority vote would report
\at{hair\_color}${=}$black on a 40/30/30 split, and would break the 40/40 tie
on \at{upper\_inner\_color} by whichever value the sort visited first.}
\label{tab:r-noise}
\end{table}

Three distinct problems appear in this track. Ties make the aggregate arbitrary.
Splits across near-synonyms (\at{short}, \at{straight}, \at{long+straight}) are
not a disagreement about the person but about the vocabulary. Occlusion
(\at{not visible} 80\%) is not a value: it is the absence of an observation, and
counting it as one would let an invisible garment outvote a visible one. A voting
aggregator cannot distinguish these cases; an LLM that is told what each form
means can.

% ========================================================= 5 the reasoning ==
\section{Multi-Agent Semantic Reasoning}
\label{sec:r-reasoning}

\subsection{The funnel}
Figure~\ref{fig:r-reasoning} shows the reasoning block. It keeps ZARA's
coarse-to-fine iteration (Section~\ref{sec:r-zara}), pruning first and deciding
second, with four agents run over two rounds:

\begin{enumerate}
\item The attribute-value selector (stage 1) chooses the attribute-value pairs to
      spend evidence on, from the knowledge base and the pool statistics.
\item The evidence pruning agent (stage 2A) reads, for every candidate and every
      selected pair, what this crop says and what the evidence around that
      reading looks like.
\item After the coarse cut ($75 \rightarrow 25$), a second selector (stage 3)
      works on the survivors, whose statistics have changed because the pool has.
\item The decision agent (stages 2B and 4B) ranks the candidates from the notes
      and returns the final ordering.
\end{enumerate}

\begin{figure}[H]
\centering
\includegraphics[width=0.95\textwidth]{reasoning-000.png}
\caption{The semantic reasoning block. The knowledge base supplies per-attribute
statistics to both selector calls; the text query reaches the pruning and
decision agents directly.}
\label{fig:r-reasoning}
\end{figure}

\paragraph{Why stage 0 is not in the funnel.} Neither the list above nor
Figure~\ref{fig:r-reasoning} shows the query analyst (stage 0), on purpose. The
funnel is the part that looks at candidates, and stage 0 never sees one. It is a
\emph{pre-run} over the query alone: it runs once per query, before the pool is
examined, and settles how each word of the query relates to each value used in
the index. Its output is a lookup table of relations, not a decision about any
candidate, and every stage of both rounds reads the same table. Drawing it as a
fifth agent inside the funnel would suggest that it takes part in pruning or
ranking, which it does not. Section~\ref{sec:r-stage0} describes it and explains
why it exists.

Two properties of this arrangement drive the prompt design. First, the coarse cut
cannot be undone: a candidate placed 26th in the first round can never be
returned, while a candidate ranked low in the second round still can. The two
mistakes are not symmetric, and the ranker is told so in round-dependent terms
(Section~\ref{sec:r-coverage}). Second, the second round renumbers the
candidates: the survivors are presented to stages 3 and 4 as candidates $1..M$
without their round-1 labels, so that the refine round is not anchored by the
ordering of the coarse round.

\subsection{Stage 0: query analyst}
\label{sec:r-stage0}
Before any candidate is seen, one call settles what the query asks of each
attribute. For every attribute in the index vocabulary it emits either
\code{NOT-CONSTRAINED} or one line per (index value, query value) pair with a
relation:

\begin{center}
\begin{tabular}{@{}ll@{}}
\toprule
\textbf{Relation} & \textbf{Meaning}\\
\midrule
\code{exact}     & interchangeable both ways; nothing gained or lost by writing either\\
\code{near}      & a plausible neighbour: what a describer might write for the same
                   thing\\
                 & under different light, from further away, or in motion\\
\code{far}       & a different value of the same kind, unlikely to be confused\\
\code{unrelated} & not the same property\\
\code{neutral}   & can neither confirm nor contradict (hides the attribute, too vague)\\
\bottomrule
\end{tabular}
\end{center}

\noindent
This is a statement about \emph{words}: the analyst sees no candidate, no image
and no statistic. The \code{near} relation matters most, because later stages use
it to tell a describer's mistake from a different person. The relations are
settled once per query and reused by every later stage, including the second
round, since asking the same question twice could produce two different answers.

\paragraph{Why a separate pre-run.} Every later stage has to decide whether what a
crop read agrees with what the query asked: the selector to weigh a pair, the
evidence reader to state where the query value stands, and the ranker to size a
conflict. Without stage 0, each of these calls would derive the relation between,
for example, \emph{navy} and \emph{black} on its own, inside a prompt already
crowded with statistics, and nothing would keep the answers consistent: the
evidence reader could treat a reading as a near match while the ranker treated
the same reading as a contradiction. Settling the relations once, from the words
alone, has three consequences.

\begin{itemize}
\item Consistency: one relation per (index value, query value) pair, read in the
      same way by stages 1 to 4 in both rounds.
\item Separation of concerns: a relation is a statement about language. Deciding
      it without candidates, images or statistics keeps it free of bias from how
      common a value is or from what the pool happens to contain; the later
      stages then reason about evidence, not vocabulary.
\item Cost: the call depends only on the query and the index vocabulary, so it is
      small, runs once per query, and is replayed from disk whenever the query is
      run again (Section~\ref{sec:r-engineering}).
\end{itemize}

\noindent
The most useful relation it produces is \code{near}: it lets the evidence reader
and the ranker treat a likely confusion of the describer (navy read as black in a
dim crop) differently from a real mismatch (navy read as red).

\subsection{Stage 1: attribute-value selector}
The selector chooses 4 to 10 attribute=value pairs with the highest expected
decision value, each with a one-sentence justification
(Figure~\ref{fig:r-selector}). It sees two tables: the \emph{global prior}, the
value's share over the whole index, and the \emph{local prior}, the pool's
distribution with the mean CLIP similarity and track stability per value
(Figure~\ref{fig:r-selector-work}).

\begin{figure}[H]
\centering
\includegraphics[width=0.90\textwidth]{selector-000.png}
\caption{The selector in the abstract. It decides \emph{where to look}, never what
is true: the output is a set of pairs with reasons, not a verdict.}
\label{fig:r-selector}
\end{figure}

\begin{figure}[H]
\centering
\includegraphics[width=0.97\textwidth]{selector_work-000.png}
\caption{The same stage with its real inputs and output: the two prior tables as
input, an attribute-value pair with its reason as output.}
\label{fig:r-selector-work}
\end{figure}

The prompt requires the choice to be made jointly across five dimensions, none of
which may decide alone: the query relation from stage 0; discrimination, that is,
how small a subset the value picks out; partitioning, how usefully it separates
competing groups; inspectability, the track stability qualified by the track
size, since a single-crop track has a stability of 1.00 by construction; and the
enrichment or depletion signal. A rare exact value can be highly decisive, a
common exact value can be weak, and a \emph{near} value can be worth a slot
because it explains an apparent mismatch.

One restriction was added after measurement: the selector's input is limited to
the attributes the query actually constrains. When shown the whole index, the
selector spent slots on attributes the query never mentioned, for which any
reading is equally consistent with the request.

\subsection{Stage 2A: evidence reader}
This is the most expensive stage and the core of the system. For each candidate
and each selected pair it writes exactly two lines: what \emph{this crop} read,
and a note about the evidence around that reading.

\begin{figure}[H]
\centering
\includegraphics[width=0.95\textwidth]{pruning-000.png}
\caption{The evidence and decision stages in the abstract: stage A reads the
evidence per candidate and per attribute, and stage B ranks from what it wrote.}
\label{fig:r-pruning}
\end{figure}

\begin{figure}[H]
\centering
\includegraphics[width=0.97\textwidth]{pruning_work-000.png}
\caption{The same two stages at work: one rendered evidence block, the note the
reader writes from it, and the ranking the decision agent produces. Every figure
in the note appears in the block; nothing is recomputed or invented.}
\label{fig:r-evidence}
\end{figure}

The rendered block for one candidate and one pair contains at most:

\begin{itemize}
\item the crop's own reading, with its relation to the query value and its index
      share, for example \code{short [near] [index 65\%]};
\item the track distribution, such as \code{short 8/10 = 80\% repeated} and
      \code{long 2/10 = 20\% repeated}, kept whole as readings rather than split;
\item linkage and cohesion: how strongly the frames that read this value link to
      this crop (\code{vs crop 0.57 +/- 0.07}) and to each other
      (\code{together 0.73 +/- 0.14}), with the number of pairs the figure is
      computed from;
\item conditions: the quality figures for the relevant band, with the deviation
      from the pool's own distribution, for example
      \code{sharp 21.5 (-1.5) mid}, \code{contrast 120.7 (+14.7) mid},
      \code{bright 88.3 (+16.2) mid}, \code{sat 110.9 (-48.1) mid};
\item the view: whether this frame is unusual for its own value, printed only
      when at least two metrics fall outside that value's own range;
\item frame agreement: how much a quoted frame resembles the crop overall,
      \emph{excluding the attribute being judged}, so that a frame agreeing about
      the upper colour cannot be used to justify trusting it about the upper
      colour.
\end{itemize}

Four rules govern the note, each adopted after a measured failure:

\begin{enumerate}
\item The citation rule: never name a mechanism without the figure that shows it.
      Adding this rule reduced the invented causes in this stage from 19 to 0,
      and the mechanisms named without numbers to 0 of 82. It is the change with
      the clearest result in this pipeline, and it was also applied to the
      ranker.
\item The crop is a reference observation, not the truth. A frame agreeing with
      the crop is not evidence that the crop is right, because a tracker can group
      two people.
\item No deciding, ranking or averaging. The reader may not write
      \emph{correct}, \emph{winner}, \emph{most likely} or \emph{therefore the
      person is X}. Its task is to make the evidence readable; ordering is done by
      another stage with another prompt.
\item A length budget: one sentence of at most 25 words for an uncontested unit,
      and 40 to 95 words where the evidence is contested. Contested notes must be
      written in full: two to four channels with their figures, then a closing
      sentence saying where the evidence leaves the readings.
\end{enumerate}

Because a single call with 75 candidates would exceed the token window of the
free tier, stage 2A is split into three calls of 25 candidates. The split is sound
because stage 2A is independent across candidates and attributes: no note depends
on another. The batches are paced (Section~\ref{sec:r-engineering}).

\subsection{Stage 2B: the ranker}
The ranker sees the query and one note per attribute per candidate, and nothing
else: no per-attribute verdict, no score and no CLIP similarity. It ranks all $N$
candidates against each other and writes at least the top
$N_{\text{keep}} = \text{cap} + 5$ rows (30 of 75 in the coarse round, 15 of 25
in the refine round) as a markdown table with one reason per row.
$N_{\text{keep}}$ is a minimum, not a quota: the model is told to keep writing
while the evidence still separates one candidate from the next, and to stop when
it no longer does.

Three mechanisms shape the ordering.

\paragraph{Reliability, not direction.} Whether an attribute matches is the first
question; how much weight the reading can bear is the second, and the ordering is
based on the second. Ranking on direction produced 69 of 75 rows opening with
``Exact red top, shorts, female gender'', all of which was shared by the whole
pool and ordered nothing.

\paragraph{What each query value is worth.} The ranker receives a rarity table
once, at the top of its message: for every query value, its share in the index,
the \emph{typical} (median) share for that attribute, the most common share, and
the track stability. Each row ends with a label, \emph{decisive}
($\geq 5\times$ typical), \emph{moderate} ($\geq 1.5\times$) or \emph{almost
nothing}, and the ratio is printed next to it so that the label can be checked.
This was added after a measured failure: a candidate whose upper body was read as
green in every frame, green being rare, was ranked below candidates that only
matched a common colour, because no downstream stage carried any notion of
rarity.

\paragraph{The coverage floor.}
\label{sec:r-coverage}
The two rounds receive different instructions about absence, because their cuts
differ in kind.

\begin{itemize}
\item In the coarse round, PRESENT means read \emph{anywhere}: in one frame, as a
      minority reading, overshadowed by a competing value, in the track but not in
      this crop, or as a \code{[near]} value standing in for the query value. A
      present but weak reading ranks above an absent one. ABSENT is an evidence
      gap in either of its two forms: \emph{read nowhere in the track} (silence)
      and \emph{nothing closer than \code{[far]} read anywhere} (contradiction).
      The floor applies in one case only, the absence of a value that the rarity
      table calls \emph{decisive}. It is neither a count of matching attributes
      nor a score.
\item In the final round nothing lies below the cut any more, so there is nothing
      left to protect. A value read nowhere is still a reliable mismatch and
      should still push a candidate down; a single overshadowed frame is weak
      positive evidence, weighed against everything else, and not a floor.
\end{itemize}

\subsection{What the CLIP score is worth inside the pool}
\label{sec:r-clipscore}
The ranker receives the evidence notes and the query, and deliberately does
\emph{not} receive the CLIP similarity. This design decision is based on the
measurement in Table~\ref{tab:r-clipscore}.

CLIP does two jobs here, and only the second is in question. Building the pool is
not contested: the 75 candidates are CLIP's top 75 out of the whole index, and no
other part of the system could have produced them. Ordering inside the pool is
what stages 2B and 4B would inherit if the score were passed on, but the pool is
already the narrow top tail of a similarity distribution, so the spread that
ranked it against the \emph{index} may be almost gone \emph{within} the pool.

\begin{table}[H]
\centering
\begin{tabular}{@{}lr@{}}
\toprule
\textbf{Measured over the evaluated pools} & \textbf{Value}\\
\midrule
Mean pool similarity                        & 0.3272\\
Mean within-pool spread (rank 1 $-$ rank 75)& 0.0622\\
Mean within-pool standard deviation         & 0.0142\\
Median adjacent-rank gap                    & 0.0004 \ (2.7\% of one sd)\\
AUC, ground truth vs.\ rest of pool         & 0.601\\
Ground-truth $z$-score within pool          & $+1.24$ sd\\
\bottomrule
\end{tabular}
\caption{What the CLIP score separates \emph{inside} the 75-candidate pool. The AUC
is the probability that a ground-truth crop outranks a non-ground-truth crop there;
1.0 would mean the score alone sorts the pool perfectly, and 0.5 would be chance.
It is computed with ties taken into account, because many crops share the same
similarity to three decimals and counting a tie as a win would overstate the score
exactly where it is weakest.}
\label{tab:r-clipscore}
\end{table}

The three measurements point the same way. An AUC of 0.601 is weak separation:
inside the pool, the score is closer to chance than to a useful ordering. The
median gap between \emph{adjacent} ranks is 2.7\% of one standard deviation, so
ranks 8 and 20 are separated by numerical noise, and a ranker told to respect them
would be respecting noise. The ground truth does sit $+1.24$ standard deviations
above the pool mean, which is above average but not separated.

CLIP is therefore a good \emph{retriever}, since it produced this pool out of a
496-identity gallery, but a weak \emph{ranker of what it retrieved}. These are
different jobs, which is why they are measured separately here. Passing the score
into the ranking prompt would add little ordering information and would anchor the
ranker to the ordering it is meant to correct.

\subsection{Stages 3 and 4}
Stage 3 is the stage-1 prompt with one paragraph inserted: it selects
attribute-value pairs again over the 25 survivors, whose pool statistics differ
from those of the original 75. A value that was rare in the pool may now be shared
by all the survivors, and therefore useless for separating them. Stage 4A is the
evidence reader again, on the new pairs, and stage 4B is the ranker with the
\code{COVERAGE\_FINAL} section. The output of 4B, cut to 10, is the system's
answer.

\subsection{Engineering}
\label{sec:r-engineering}
Several implementation details are essential for reproducibility:

\begin{itemize}
\item Replay by byte-identical input. Every stage writes its rendered input and
      its reply to disk. Running an identity again replays from disk every stage
      whose input text is byte-identical, and resumes at the first stage without a
      saved answer. A run interrupted by a quota limit can therefore continue
      without repeating the calls already made.
\item Pacing. A rolling token-per-minute pacer holds back a call that would exceed
      the window, with a fixed 60-second gap between 2A batches, 180 seconds
      between identities, and exponential retry (60/120/240/480\,s) on errors 503
      and 429.
\item Robust parsing. A duplicate row in a ranking and a short reply are different
      failures and are reported separately; a candidate skipped in the middle of
      the list is appended in input order rather than dropped; and a ranking that
      cannot be parsed falls back to the CLIP order with a warning, because an
      empty round would score exactly like a wrong answer while hiding the parse
      failure.
\end{itemize}

% ================================================================= 6 setup ==
